% =====================================================================
%  CUESTIONARIO FINAL DE CONTEXTO (al terminar la secuencia)
%  No mide aprendizaje: da contexto para interpretar lo observado.
%  El bloque B repite los ítems 6-12 del cuestionario inicial.
%  Compilar con: pdflatex cuestionario_final.tex (dos veces)
% =====================================================================
\documentclass[11pt,letterpaper]{article}

\usepackage[utf8]{inputenc}
\usepackage[T1]{fontenc}
\usepackage[spanish,es-tabla]{babel}
\usepackage{lmodern}
\usepackage{microtype}
\usepackage[margin=2cm]{geometry}
\usepackage{amssymb}
\usepackage{tabularx,array}
\usepackage{enumitem}
\usepackage[table]{xcolor}
\usepackage[most]{tcolorbox}
\usepackage{fancyhdr}
\usepackage{lastpage}
\usepackage{pgffor}

\definecolor{unalverde}{RGB}{0,98,65}
\definecolor{grisclaro}{RGB}{242,242,242}

\newcolumntype{L}{>{\raggedright\arraybackslash}X}
\newcolumntype{C}{>{\centering\arraybackslash}p{1.25cm}}
\setlist{itemsep=2pt,topsep=3pt}
\setlength{\parindent}{0pt}

\newcommand{\cas}{$\square$\,}
\newcommand{\cabeza}[2]{\footnotesize\bfseries\shortstack{#1\\#2}}
\newcommand{\bloque}[1]{\par\medskip{\large\bfseries #1}\par\smallskip}
\newcommand{\linea}{\par\rule{\linewidth}{0.4pt}\par\vspace{4pt}}
\newcommand{\lineas}[1]{\par\vspace{2pt}\foreach\i in {1,...,#1}{\linea}}
\newcommand{\escala}{& \cas & \cas & \cas & \cas \\ \hline}

\pagestyle{fancy}
\fancyhf{}
\fancyhead[L]{\small Cuestionario final --- Tutoría personalizada GEA}
\fancyhead[R]{\small Código: \rule{1.6cm}{0.4pt} \quad Fecha: \rule{1.8cm}{0.4pt}}
\fancyfoot[C]{\small \thepage\ de \pageref{LastPage}}
\renewcommand{\headrulewidth}{0.4pt}

\begin{document}

{\Large\bfseries Cuestionario final\par}
\medskip

{\itshape
\small
Ya terminó las tareas de la tutoría. Este cuestionario no es un examen y no afecta sus notas. Use el mismo código del cuestionario inicial y no escriba su nombre. Nos interesa su opinión sincera, también si algo no le sirvió. Tiempo aproximado: 10 minutos.
\par}

\bloque{A. Cuando resuelvo problemas de matemáticas}

Marque qué tan de acuerdo está \emph{hoy} con cada frase.

\begin{center}\small
\renewcommand{\arraystretch}{1.45}
\begin{tabularx}{\linewidth}{|L|C|C|C|C|}
\hline
 & \cabeza{Nada de}{acuerdo} & \cabeza{Poco de}{acuerdo} & \cabeza{De}{acuerdo} & \cabeza{Muy de}{acuerdo} \\ \hline
1. Entiendo lo que se explica en clase. \escala
2. Puedo resolver ejercicios parecidos a los que se hicieron en clase. \escala
3. Puedo resolver un problema que no se parece a ninguno de los de clase. \escala
4. Cuando resuelvo un ejercicio, puedo explicar por qué cada paso es válido. \escala
5. Cuando me equivoco, puedo encontrar dónde estuvo el error. \escala
6. Antes de dar una respuesta por buena, la compruebo de alguna manera. \escala
7. Si no veo cómo empezar un problema, intento algo de todos modos. \escala
\end{tabularx}
\end{center}

\bloque{B. Sobre las tareas y sus tres momentos}

\begin{center}\small
\renewcommand{\arraystretch}{1.45}
\begin{tabularx}{\linewidth}{|L|C|C|C|C|}
\hline
 & \cabeza{Nada de}{acuerdo} & \cabeza{Poco de}{acuerdo} & \cabeza{De}{acuerdo} & \cabeza{Muy de}{acuerdo} \\ \hline
8. Trabajar primero sin IA me sirvió para saber qué entendía y qué no. \escala
9. Con Gemini encontré errores que yo solo no había visto. \escala
10. Gemini me dio respuestas que copié sin entenderlas del todo. \escala
11. Habría preferido que Gemini me diera la respuesta directamente. \escala
12. Comparar lo azul con lo rojo me ayudó a decidir mi respuesta final. \escala
13. Las tareas me obligaron a pensar más que los ejercicios de clase. \escala
\end{tabularx}
\end{center}

\textbf{14.} ¿Cuál de las tareas le sirvió más y por qué?
\lineas{2}

\textbf{15.} ¿En qué momento le fue más útil Gemini? ¿Y en qué momento le estorbó o lo confundió?
\lineas{3}

\textbf{16.} Piense en una de las tareas. ¿Qué parte de la solución hizo usted y qué parte hizo Gemini?
\lineas{3}

\bloque{C. Lo que viene}

\textbf{17.} La próxima vez que estudie solo y se quede atascado, ¿qué va a hacer?
\lineas{2}

\textbf{18.} ¿Va a seguir intentando primero por su cuenta antes de usar una IA? \quad \cas Sí \quad \cas A veces \quad \cas No \quad ¿Por qué?
\lineas{2}

\textbf{19.} ¿Qué le cambiaría a esta forma de trabajar en la tutoría?
\lineas{2}

\vspace{4pt}
{\small\centering Gracias por participar.\par}

\end{document}
